% !PRINTED 770
\weberpage{770}

Ist $p$ eine in $n$ aufgehende Primzahl, so ist
\weq{3}{\frac{x^n-1}{x^{\frac{n}{p}}-1} = x^{\frac{n}{p}(p-1)} + x^{\frac{n}{p}(p-2)} + \cdots + x^{\frac{n}{p}} + 1,}
und diese Function verschwindet für $x=\alpha, \beta, \gamma, \dots$ und ist
daher durch $f_n(x)$ theilbar. Wir wollen sie $= f_n(x)\, g(x)$ setzen,
worin dann $g(x)$ (nach Bd. I, \S.\,2) eine ganze Function von $x$
mit ganzzahligen Coëfficienten ist.

Ist nun
\[
n = p^k n',
\]
$n'$ nicht mehr durch $p$ theilbar, und $\alpha'$ eine primitive $n'^{\mathrm{te}}$ Einheitswurzel, so ist
\[
\alpha'^{\frac{n}{p}} = 1,
\]
und folglich, wenn man $x = \alpha'$ setzt [nach (3)]:
\weq{4}{f_n(\alpha')\, g(\alpha') = p.}

Wir leiten hieraus zunächst einen Beweis für die Irreducibilität von $f_n(x)$ unter der Voraussetzung ab, dass $n$ eine
\gesperrt{Primzahlpotenz} ist, wenn auch für diesen Fall der in Bd. I,
\S.\,134 gegebene Beweis einwandsfrei ist.

Angenommen, es zerfalle $f_n(x)$ in zwei rationale Factoren:
\weq{5}{f_n(x) = \varphi(x)\, \psi(x).}

Dann sind alle Wurzeln, sowohl von $\varphi(x)$ als von $\psi(x)$,
$n^{\mathrm{te}}$ Einheitswurzeln, und ihre $n^{\mathrm{ten}}$ Potenzen also gleich 1. Wenn
nun $n = p^k$ ist, so können wir aus $\varphi(x)$ und $\psi(x)$ zwei ganze
Functionen $\Phi(x)$, $\Psi(x)$ ableiten, deren Wurzeln die $n^{\mathrm{ten}}$ Potenzen der Wurzeln von $\varphi(x)$ und $\psi(x)$ sind und die daher alle
gleich 1 sind, und aus (2) ergiebt sich:
\[
\varphi(x) \equiv \Phi(x), \quad \psi(x) \equiv \Psi(x) \ (\mathrm{mod}\ p),
\]
woraus für $x = 1$ folgt:
\[
\varphi(1) \equiv 0, \quad \psi(1) \equiv 0 \ (\mathrm{mod}\ p).
\]
Hiernach ergiebt sich aus (5):
\weq{6}{f_n(1) \equiv 0 \ (\mathrm{mod}\ p^2).}

Nun ist aber in diesem Falle $f_n(1) = p$ (Bd. I, \S.\,133, V.),
was der Formel (6) widerspricht. Unsere Annahme war also
unzulässig und $f_n(x)$ ist irreducibel.
